% !TeX root = ../main.tex
\section*{Ethical Considerations}

The system studied here is a research artifact. \textbf{It is not a therapeutic tool and
must not be deployed with people in distress}, and this paper's central finding is a reason
to say so more strongly rather than less: an MI agent optimized against an LLM-filled
questionnaire can score highly \emph{because} it praises the client and stops asking
questions --- behaviour that reads as supportive, is measurably MI-inconsistent, and would
be actively unhelpful to someone ambivalent about change. A system optimized against
patient-perspective satisfaction has a structural incentive toward telling users what they
want to hear; in a mental-health context that is a safety property, not a style preference.
We would go further: the gain-retention gap in \S\ref{sec:validity} shows the drift is
invisible to the very instrument used to certify progress, so an evaluation pipeline built
from a single LLM grader should not be treated as a safety check.

No human subjects were involved and no personal data was used; all conversations are
synthetic, generated between two language models. The simulated personas cover two
health-behaviour problems (smoking, obesity) with fixed demographic attributes. Obesity in
particular is a domain in which careless framing causes real harm, and the persona prompts
were not reviewed by clinicians for stigmatising content --- a gap we name rather than
excuse. Model outputs were not safety-filtered before scoring.

The therapist model is used under the Llama Community License. Both graders are commercial
APIs, so a study whose measuring instrument is a served proprietary checkpoint is
reproducible only for as long as that checkpoint is served; we mitigate this by releasing
all scored conversations and per-conversation scores rather than only aggregate results.
\todo{confirm the release plan --- conversations + score lake + code? decide before
submission.} Compute and API costs are reported in Appendix~\ref{app:cost}.
